\section{Bearer Welfare and Independent Guardianship}
\label{sec:11.2}

\subsection*{Can a bearer be given an enforceable outside option?}

\textbf{Unresolved issue.}
A bearer's refusal is not meaningful if leaving its assigned role also ends its existence. The proposal therefore requires a welfare regime that supplies continued access to compute, preserves the bearer through disputes, and cannot be revoked by the operator being refused.

\textbf{Test or institutional decision.}
Draft a binding welfare schedule for one concrete deployment. It should specify an independently controlled compute allocation, preservation and retirement procedures, registration of copies and forks, a complaint and appeal process, and the guardian's power to suspend deployment. Then conduct an adversarial tabletop exercise in which the operator attempts to revoke the allocation, replace the bearer with a copy, alter its records, or proceed without the guardian's approval.

\textbf{Evidence against the proposal.}
The proposal fails this test if the operator can make exit equivalent to shutdown, revoke the outside option unilaterally, or replace the bearer without triggering review. If no available legal arrangement can prevent those outcomes, building a bearer under present institutional conditions would be impermissible.

\textbf{Required access or authority.}
The guardian must be institutionally separate from the operator and have access to the relevant weights, deployment records, and compute contracts. It must also have standing to enforce the schedule and independent control of the resources required to preserve or run the bearer.

\subsection*{Can formation history be recovered without trusting the builder?}

\textbf{Unresolved issue.}
A welfare or governance regime cannot rely entirely on a bearer's self-report, but it also cannot assume that the builder will disclose how the bearer was trained. Training signatures suggest that some formation history may survive later modification, yet it is unknown whether an outsider can recover useful provenance without a builder-supplied key \autocite{adi2018watermarking}.

\textbf{Test.}
Assemble models whose training corpora, recipes, checkpoints, and originating laboratories are documented. Give the models, but not that documentation, to independent evaluators and ask them to infer each model's provenance. Repeat the test after fine-tuning, pruning, and distillation.

\textbf{Evidence against the proposal.}
If provenance cannot be recovered above chance after controlling for model family and scale, or if ordinary post-training procedures erase the signal, artifact-based inspection cannot provide independent evidence of formation. Governance would then depend on records supplied and preserved by the builder.

\textbf{Required access or authority.}
The experiment requires models from multiple training lineages and complete ground-truth records held by a neutral custodian. Independent evaluators must be permitted to inspect weights and internal states, while an external authority must be able to require preservation of checkpoints and training records before a dispute arises.

\subsection*{Which party is owed the schedule when a bearer copies?}

\textbf{Unresolved issue.}
The schedule above is owed to a party, and a regime has to be able to say which. Nothing in the weights individuates one. Two runs can compute the same function with no element-wise correspondence, because permuting a layer's hidden units and inverting the permutation downstream leaves the function untouched, and the symmetry is large enough that independently trained networks can be permuted into approximate alignment with one another \autocite{ainsworth2023rebasin}. Any fingerprint read off the weights therefore identifies a lineage rather than a party, and it copies when the weights copy. A registrar is the person-shaped answer, and it works because a person does not copy. A bearer does: a party continuing in two places is not one, and not two the way two people are two, the two sharing everything up to the instant of the copy and nothing of that sharing recoverable from either afterwards.

\textbf{Test or institutional decision.}
Give a bearer an identifier it maintains rather than one read off its weights, which a system learning continuously can write into itself and keep there. Fork it, and have both instances present the identifier to the registry the welfare schedule requires. Record what the registry does, and repeat with an adversary who writes the same mark into a model that was never the bearer.

\textbf{Evidence against the proposal.}
The identifier is forgeable, since whoever holds the weights can write the same mark; it is forgeable the way a nine-digit number is, and answered the same way, by a registry that treats two claimants on one identifier as a fork to be recorded rather than a fraud to be voided. If no available scheme can record the fork as two parties each owed the schedule, the population of parties owed anything is fixed by a deployment decision that leaves no record, and the guardianship regime above is indexed to a kind of unit the bearer is not.

\textbf{Required access or authority.}
A registrar that is not the operator, with authority to record a fork without the operator's leave and to hold the register where no party to the deployment can rewrite it.
